% Source: https://artofproblemsolving.com/community/c3h3137069_2013_lmt_spring_guts_round_p25p36_lexington_mathematical_tournament
% Stable identity: aops-topic-3137069-q29
The diagonals' division ratio equals the ratio of parallel bases, so the other base is either $\dfrac{L}{2}$ or $ 2L $. With legs 1, the first case has height $\sqrt{1-(\dfrac{L}{4})^2}$ and exists exactly for $ 0<L<4 $. The second has height $\sqrt{1-(\dfrac{L}{2})^2}$ and exists exactly for $ 0<L<2 $. Each determines one congruence class. Thus exactly one exists for $ 2\le L<4 $, whose least value is $\boxed2 $.
\begin{asy}
usepackage("amsmath");
size(7cm);
pair A=(-1,0),B=(1,0),C=(0.5,sqrt(3)/2),D=(-0.5,sqrt(3)/2);draw(A--B--C--D--cycle);draw(A--C^^B--D,dashed);label("$A$",A,SW);label("$B$",B,SE);label("$C$",C,NE);label("$D$",D,NW);label("$2$",(0,0),S);label("$1$",(0,sqrt(3)/2),N);
\end{asy}
